\centerline{\bf \Huge Принцип Дирихле II}

\begin{flushright}
        \scriptsize{ДИРИХЛЕЕЕЕЕЕЕЕЕЕ! - Алексей Владимирович Савватеев}
\end{flushright}


\problem{\textbf{1}} Олимпиаду писали 70 школьников. Аркаша набрал 33 балла, остальные меньше. Докажите, что по крайней мере три школьника набрали одинаковое количество баллов.





\problem{\textbf{2}} В кинотеатре 7 рядов по 10 мест. Группа из 50 человек сходила на утренний и вечерний сеансы. Докажите, что найдутся два человека, которые и утром, и вечером сидели на одном ряду.



\problem{\textbf{3}} Можно ли таблицу 2025×2025 заполнить числами -1, 0, 1 так, чтобы суммы во всех строках, во всех столбцах и на главных диагоналях были различны?

\begin{flushright} \textbf{2 балла}\\ 
\end{flushright}


\problem{\textbf{4}} Докажите, что если $a, b, c$ — нечётные числа, то хотя бы одно из чисел $ab - 1, bc - 1, ca - 1$ делится на 4.



\problem{\textbf{5}} 
По кругу записаны 2025 натуральных чисел. Известно, что в каждой паре соседних чисел одно делится на другое. Докажите, что найдётся пара несоседних чисел с таким же свойством.



\problem{\textbf{6}}  В классе 25 человек. Известно, что среди любых трех из них есть двое друзей. Докажите, что есть ученик, у которого не менее 12 друзей.\




\problem{\textbf{7}} Докажите, что найдется число вида 777…77000..00, делящееся на 2025.


\problem{\textbf{8}} В воздушном пространстве находятся облака. Оказалось, что пространство можно разбить десятью плоскостями на части так, чтобы в каждой из частей находилось не более одного облака. Через какое наибольшее число облаков мог пролететь самолет, придерживаясь прямолинейного курса?

\begin{flushright} \textbf{3 балла}\\ 
\end{flushright}